\documentclass[11pt]{article}

\usepackage[a4paper,margin=1in]{geometry}
\usepackage{amsmath,amssymb,amsthm,mathtools}
\usepackage{bm}
\usepackage{graphicx}
\usepackage{float}
\usepackage{booktabs}
\usepackage{xcolor}
\usepackage{hyperref}
\usepackage{listings}
\usepackage{enumitem}
\usepackage{caption}
\usepackage{subcaption}
\usepackage{microtype}

\graphicspath{{hw1_solution_figs/}}

\hypersetup{
  colorlinks=true,
  linkcolor=blue,
  urlcolor=blue,
  citecolor=blue
}

\lstdefinestyle{pythonstyle}{
  language=Python,
  basicstyle=\ttfamily\footnotesize,
  keywordstyle=\color{blue!70!black},
  commentstyle=\color{green!40!black},
  stringstyle=\color{red!60!black},
  showstringspaces=false,
  breaklines=true,
  frame=single,
  columns=fullflexible,
  keepspaces=true,
  tabsize=2
}

\newtheorem*{claim}{Claim}
\newtheorem*{lemma}{Lemma}
\newcommand{\R}{\mathbb{R}}
\newcommand{\relu}{\operatorname{ReLU}}
\newcommand{\norm}[1]{\left\lVert #1\right\rVert}
\newcommand{\ip}[2]{\left\langle #1,#2\right\rangle}

\title{Mathematical Theory for Deep Neural Networks\\HW1 Solution}
\author{Wells Guan\\University of Wisconsin--Madison}
\date{July 6, 2026}

\begin{document}
\maketitle

\section*{Problem 1. Energy error estimate for gradient descent}

Let $F:\R^d\to\R$ be continuously differentiable, strictly convex, and assume that $\nabla F$ is $L$-Lipschitz. Let $u$ be a minimizer of $F$, and define
\[
  u^m=u^{m-1}-\tau \nabla F(u^{m-1}),\qquad 0<\tau\le \frac1L.
\]

First, $u$ is unique. Indeed, if $u$ and $v$ are two different minimizers, then for any $t\in(0,1)$ strict convexity gives
\[
  F(tu+(1-t)v)<tF(u)+(1-t)F(v)=\min F,
\]
which is impossible. Hence the minimizer is unique. Also, since the minimization is over all of $\R^d$ and $F$ is differentiable,
\[
  \nabla F(u)=0.
\]

We now prove the estimate. The $L$-smoothness of $F$ gives the descent lemma
\[
  F(y)\le F(x)+\ip{\nabla F(x)}{y-x}+\frac L2\norm{y-x}^2.
\]
Taking $y=x-\tau\nabla F(x)$, we obtain
\[
  F(x-\tau\nabla F(x))
  \le F(x)-\tau\left(1-\frac{L\tau}{2}\right)\norm{\nabla F(x)}^2
  \le F(x)-\frac{\tau}{2}\norm{\nabla F(x)}^2.
\]
In particular, the sequence $F(u^m)$ is nonincreasing.

Set $x=u^{m-1}$ and $x^+=u^m=x-\tau\nabla F(x)$. By convexity,
\[
  F(x)-F(u)\le \ip{\nabla F(x)}{x-u}.
\]
Furthermore,
\[
\begin{aligned}
  \norm{x-u}^2-\norm{x^+-u}^2
  &=2\tau\ip{\nabla F(x)}{x-u}-\tau^2\norm{\nabla F(x)}^2 \\
  &\ge 2\tau(F(x)-F(u))-\tau^2\norm{\nabla F(x)}^2.
\end{aligned}
\]
From the descent inequality,
\[
  \tau^2\norm{\nabla F(x)}^2
  \le 2\tau\bigl(F(x)-F(x^+)\bigr).
\]
Therefore,
\[
\begin{aligned}
  \norm{x-u}^2-\norm{x^+-u}^2
  &\ge 2\tau(F(x)-F(u))-2\tau(F(x)-F(x^+)) \\
  &=2\tau(F(x^+)-F(u)).
\end{aligned}
\]
Thus, for every $k\ge1$,
\[
  2\tau\bigl(F(u^k)-F(u)\bigr)
  \le \norm{u^{k-1}-u}^2-\norm{u^k-u}^2.
\]
Summing from $k=1$ to $m$ gives
\[
  2\tau\sum_{k=1}^m \bigl(F(u^k)-F(u)\bigr)
  \le \norm{u^0-u}^2-\norm{u^m-u}^2
  \le \norm{u^0-u}^2.
\]
Since $F(u^k)$ is nonincreasing,
\[
  m\bigl(F(u^m)-F(u)\bigr)
  \le \sum_{k=1}^m \bigl(F(u^k)-F(u)\bigr).
\]
Hence
\[
  \boxed{
  F(u^m)-F(u)
  \le \frac{\norm{u^0-u}^2}{2\tau m}
  }
\]
for every $m\ge1$.

\section*{Problem 2. Linear regions of random ReLU networks}

For the numerical experiment, I used fully connected ReLU networks with input dimension $2$. For $L$ hidden layers, set
\[
  x^0=x\in\R^2,
  \qquad
  x^\ell=\relu(W^{\ell-1}x^{\ell-1}+b^{\ell-1}),\quad \ell=1,\dots,L,
\]
and take the scalar output
\[
  u(x)=W^Lx^L+b^L.
\]
The width is $n_\ell=10$ for each hidden layer. The parameters are sampled independently from the standard normal distribution. On each region where all signs of the pre-activations are fixed, the ReLU network is affine; hence these sign patterns determine the linear regions.

The following figures were generated by sampling a fine grid on $[-1,1]^2$, computing the binary activation pattern at each point, and assigning one color to each distinct pattern.

\begin{figure}[H]
\centering
\begin{subfigure}{0.48\textwidth}
  \centering
  \includegraphics[width=\textwidth]{problem2_L1_regions.png}
  \caption{$L=1$, one hidden layer of width $10$.}
\end{subfigure}
\hfill
\begin{subfigure}{0.48\textwidth}
  \centering
  \includegraphics[width=\textwidth]{problem2_L2_regions.png}
  \caption{$L=2$, two hidden layers, each of width $10$.}
\end{subfigure}
\caption{Sampled linear regions of random ReLU networks on $[-1,1]^2$.}
\end{figure}

\section*{Problem 3. Relationship between the hat and ReLU stiffness matrices}

Let
\[
  0=x_0<x_1<\cdots<x_{N+1}=1,
  \qquad x_i=ih,
  \qquad h=\frac1{N+1}.
\]
Put $M=N+1$. We use the standard one-dimensional hat functions $\{\phi_i\}_{i=1}^M$, with the usual endpoint modification for $\phi_M$ at $x=1$. Thus $\phi_i(x_j)=\delta_{ij}$ for $1\le i,j\le M$, $\phi_i$ is piecewise linear, and $\phi_i(0)=0$.

The stiffness matrix for the hat basis is
\[
  (A_\phi)_{ij}=\int_0^1 \phi_i'(x)\phi_j'(x)\,dx.
\]
A direct calculation gives
\[
  A_\phi=\frac1h K,
\]
where
\[
  K=
  \begin{pmatrix}
  2 & -1 \\
  -1 & 2 & -1 \\
  & \ddots & \ddots & \ddots \\
  && -1 & 2 & -1 \\
  &&& -1 & 1
  \end{pmatrix}\in\R^{M\times M}.
\]
The last diagonal entry is $1$, because the right endpoint $x=1$ has a natural boundary condition side and the last hat function only has one element in its support.

Now define the ReLU basis
\[
  \psi_i(x)=\relu(x-x_{M-i}),\qquad i=1,\dots,M.
\]
Since
\[
  \psi_i'(x)=\mathbf 1_{(x_{M-i},1)}(x)
\]
up to values at the grid points, we have
\[
\begin{aligned}
  (A_\psi)_{ij}
  &=\int_0^1 \psi_i'(x)\psi_j'(x)\,dx \\
  &=\left| (x_{M-i},1)\cap (x_{M-j},1)\right| \\
  &=1-\max\{x_{M-i},x_{M-j}\} \\
  &=h\min\{i,j\}.
\end{aligned}
\]
Therefore
\[
  A_\psi=Q,
  \qquad
  Q_{ij}=h\min\{i,j\}.
\]
It remains to check that $Q=A_\phi^{-1}$. Since $A_\phi=(1/h)K$, it is enough to show
\[
  K B=I,
  \qquad B_{ij}=\min\{i,j\}.
\]
Fix a column $j$ and set $v_i=B_{ij}=\min\{i,j\}$. For $1<i<M$, the $i$th component of $Kv$ is
\[
  (Kv)_i=-v_{i-1}+2v_i-v_{i+1}.
\]
If $i<j$, then
\[
  -v_{i-1}+2v_i-v_{i+1}=-(i-1)+2i-(i+1)=0.
\]
If $i>j$, then
\[
  -v_{i-1}+2v_i-v_{i+1}=-j+2j-j=0.
\]
If $i=j<M$, then
\[
  -v_{j-1}+2v_j-v_{j+1}=-(j-1)+2j-j=1.
\]
At the left endpoint the same formula gives $(Kv)_1=1$ when $j=1$ and $0$ otherwise. At the last row,
\[
  (Kv)_M=-v_{M-1}+v_M,
\]
which is $0$ if $j<M$ and $1$ if $j=M$. Hence $KB=I$. Thus
\[
  K^{-1}=B,
  \qquad
  A_\phi^{-1}=hK^{-1}=Q=A_\psi.
\]
Therefore
\[
  \boxed{A_\psi=A_\phi^{-1}.}
\]

\section*{Problem 4. FEM for a one-dimensional Poisson problem and GD solution}

Consider
\[
  -u''=f\quad\text{in }(0,1),
  \qquad
  u(0)=0,
  \qquad
  u'(1)=0,
\]
where
\[
  f(x)=\frac94\pi^2\sin\left(\frac32\pi x\right),
  \qquad
  u(x)=\sin\left(\frac32\pi x\right).
\]
Indeed,
\[
  -u''(x)=\frac94\pi^2\sin\left(\frac32\pi x\right)=f(x),
\]
and
\[
  u(0)=0,
  \qquad
  u'(1)=\frac32\pi\cos\left(\frac32\pi\right)=0.
\]

The weak formulation is: find $u\in V$ with $u(0)=0$ such that
\[
  a(u,v)=(f,v),\qquad \forall v\in V,
\]
where
\[
  a(u,v)=\int_0^1 u'(x)v'(x)\,dx,
  \qquad
  (f,v)=\int_0^1 f(x)v(x)\,dx.
\]
The Neumann condition at $x=1$ is natural and is already included in the weak formulation.

For a basis $b_1,\dots,b_M$, where $b_i=\phi_i$ for the hat basis or $b_i=\psi_i$ for the ReLU basis, write
\[
  u_h(x)=\sum_{i=1}^M c_i b_i(x).
\]
Then the finite element system is
\[
  A_b c=g_b,
\]
with
\[
  (A_b)_{ij}=\int_0^1 b_i'(x)b_j'(x)\,dx,
  \qquad
  (g_b)_i=\int_0^1 f(x)b_i(x)\,dx.
\]

To solve this linear system by gradient descent, minimize the quadratic energy
\[
  J_b(c)=\frac12 c^T A_b c-g_b^T c.
\]
The gradient is
\[
  \nabla J_b(c)=A_b c-g_b.
\]
Thus the GD iteration is
\[
  c^{(m+1)}=c^{(m)}-\tau_b(A_b c^{(m)}-g_b),
  \qquad
  c^{(0)}=0,
\]
where I use
\[
  \tau_b=\frac1{\lambda_{\max}(A_b)}.
\]

For the numerical plots below, I used $N=31$, so $M=N+1=32$ unknowns. Since the exact solution is known, I plot the error, also called the residual function here,
\[
  r_b^{(m)}(x)=u_{h,b}^{(m)}(x)-u(x),
\]
for $m=1,50,500,5000$.

\begin{figure}[H]
\centering
\begin{subfigure}{0.48\textwidth}
  \centering
  \includegraphics[width=\textwidth]{problem4_hat_residuals.png}
  \caption{Hat basis.}
\end{subfigure}
\hfill
\begin{subfigure}{0.48\textwidth}
  \centering
  \includegraphics[width=\textwidth]{problem4_relu_residuals.png}
  \caption{ReLU basis.}
\end{subfigure}
\caption{Residual/error functions $r_b^{(m)}(x)=u_{h,b}^{(m)}(x)-u(x)$ at selected GD iterations.}
\end{figure}

The hat and ReLU bases span the same finite element space, but the coefficient systems are represented in different coordinates. In particular, by Problem 3, $A_\psi=A_\phi^{-1}$, so the two GD iterations use different matrices and different optimal time steps.

\section*{Problem 5. Universal approximation by one-hidden-layer networks}

The statement in the problem is usually understood in the uniform-density sense:
\[
  \overline{\Sigma^\sigma}^{\|\cdot\|_{L^\infty(\Omega)}}\supset C(\Omega).
\]
The literal inclusion $\Sigma^\sigma\supset C(\Omega)$ is generally false; for example, finite ReLU sums in one dimension are continuous piecewise linear functions, not all continuous functions. Therefore I prove the standard density statement. Assume $\Omega\subset\R^d$ is compact; if $\Omega$ is not compact, the statement should be interpreted locally on compact subsets.

Recall
\[
  \Sigma^\sigma
  =\left\{
  \sum_{i=1}^n a_i\sigma(w_i\cdot x+b_i):
  a_i,b_i\in\R,
  w_i\in\R^d,
  n\ge0
  \right\}.
\]

\begin{lemma}[Discriminatory property of non-polynomial activations]
Let $\sigma$ be Riemann integrable and not equal almost everywhere to a polynomial. If $\mu$ is a finite signed Borel measure on a compact set $\Omega\subset\R^d$ such that
\[
  \int_\Omega \sigma(w\cdot x+b)\,d\mu(x)=0,
  \qquad \forall w\in\R^d,\ \forall b\in\R,
\]
then $\mu=0$.
\end{lemma}

\begin{proof}[Sketch of the lemma]
The key idea is that a non-polynomial activation has enough ridge translates and dilates to test all Fourier modes. If necessary, first mollify $\sigma$ by a smooth compactly supported mollifier. The annihilation identity is preserved under mollification because
\[
  (\sigma*\rho)(w\cdot x+b)
  =\int_\R \sigma(w\cdot x+b-s)\rho(s)\,ds.
\]
Since $\sigma$ is not a polynomial a.e., one can choose the mollifier so that $\sigma*\rho$ is still non-polynomial. For a fixed $w$, push $\mu$ forward by $x\mapsto w\cdot x$. The above identity says that the convolution of this push-forward measure with $\sigma*\rho$ is zero. Taking the one-dimensional Fourier transform implies that the Fourier transform of the push-forward measure vanishes at every frequency where $\widehat{\sigma*\rho}$ is nonzero. Since $\sigma*\rho$ is not a polynomial, its Fourier transform is not supported only at the origin. Dilating $w$ then gives
\[
  \widehat{\mu}(\xi)=\int_\Omega e^{-i\xi\cdot x}\,d\mu(x)=0,
  \qquad \forall \xi\in\R^d.
\]
The Fourier transform uniquely determines a finite measure on a compact set, so $\mu=0$.
\end{proof}

Now let
\[
  X=\overline{\Sigma^\sigma}^{\|\cdot\|_\infty}\subset C(\Omega).
\]
Suppose, for contradiction, that $X\ne C(\Omega)$. By the Hahn--Banach theorem and the Riesz representation theorem, there exists a nonzero finite signed Borel measure $\mu$ on $\Omega$ such that
\[
  \int_\Omega v(x)\,d\mu(x)=0,
  \qquad \forall v\in X.
\]
In particular, since every ridge function $x\mapsto \sigma(w\cdot x+b)$ lies in $\Sigma^\sigma\subset X$, we have
\[
  \int_\Omega \sigma(w\cdot x+b)\,d\mu(x)=0,
  \qquad \forall w\in\R^d,\ \forall b\in\R.
\]
By the discriminatory lemma, $\mu=0$, contradicting the choice of $\mu$. Hence
\[
  X=C(\Omega),
\]
or equivalently,
\[
  \boxed{
  \overline{\Sigma^\sigma}^{\|\cdot\|_\infty}=C(\Omega).
  }
\]
Thus one-hidden-layer neural networks with a non-polynomial activation are dense in $C(\Omega)$.

\appendix
\section*{Appendix. Python code used for Problems 2 and 4}

\lstinputlisting[style=pythonstyle]{hw1_algorithms.py}

\end{document}
